\chapter{Canonical comparison on the original relative family}

\begin{theorem}[Actual fibre cohomology from free localized homology]
\label{thm:relative-native-fibre-free-specialization}
\uses{thm:relative-arithmetic-free-scalar-homology,lem:kn-derived-geometric-fibre-cohomology}
\lean{EnvelopeProperty.Vanishing.RelativeCanonicalSpecialization.fibre_subsingleton_of_free_scalar_fibre}
\leanok
Let $Y/A$ be separated and let $M$ be invertible on $Y$. Fix a finite
ordered affine cover whose Čech complex $K$ is bounded and has flat
$A$-module terms. Let $A\to T\to C$ be a compatible scalar tower with
$C\ne0$, and suppose that the native homology of $T\otimes_A K$ is
free over $T$ in every degree. If $H^{q+1}(Y_C,M_C)=0$, then
$H^{q+1}(Y_F,M_F)=0$ for every compatible $T$-algebra $F$.
\end{theorem}
\begin{proof}
The actual inverse-image affine cover computes cohomology of each scalar
fibre, and its ordered complex is canonically the corresponding tensor
complex. Below the upper bound, the preceding homology base-change theorem
identifies $H^{q+1}(C\otimes_T(T\otimes_A K))$ with
$C\otimes_T H^{q+1}(T\otimes_A K)$. Associativity of tensor products
identifies the former with the stipulated zero group. A free module whose
extension to a nonzero ring is zero has empty basis, so the original
$T$-homology vanishes. Its extension to $F$ is zero, and the same homology
base-change theorem gives the desired fibre vanishing. Above the bound,
the actual middle Čech term and its scalar extensions are zero.
\end{proof}

\begin{theorem}[Canonical specialization at the produced Witt point]
\label{thm:relative-canonical-produced-witt-specialization}
\uses{thm:relative-native-fibre-free-specialization,lem:relative-arithmetic-coordinate-complex,lem:relative-arithmetic-cech-free-witt-point,lem:kn-arithmetic-canonical-specialization-model}
\lean{EnvelopeProperty.Vanishing.RelativeCanonicalSpecializationWitt.nativeCanonicalTensor_subsingleton}
\leanok
Let $A$ be a characteristic-zero noetherian domain smooth and of finite
type over $\mathbf Z$, and let $B$ be a noetherian finite-type $A$-algebra.
Let $Y/A$ be smooth and separated of relative dimension $n$, with a
compatible projective map $Y\to\operatorname{Spec}B$, and let $L$ be
invertible on $Y$. Fix $q\geq0$. Suppose that every finite quotient field
$A/P$ of characteristic $\ell>n$ admitting an actual lift
$A\to W_2(A/P)$ has
$H^{q+1}(Y_{A/P},\omega_{Y_{A/P}/(A/P)}\otimes L_{A/P})=0$.
Then the same canonical-twisted group vanishes on $Y_F$ for every
injective homomorphism $A\to F$ into a field.
\end{theorem}
\begin{proof}
Construct $M=\omega_{Y/A}\otimes L$ and its invertibility. Smoothness makes
$Y/A$ flat. Relative projectivity constructs the actual bounded Čech
complex, with flat $A$-module terms and homology finite over $B$.
The preceding free-Witt-point lemma chooses an actual $P$ and
$\ell>n$, its finite quotient field and Witt lift, making the localized
complex's homology free. The stipulated geometric vanishing applies at
this chosen point. Canonical base change identifies its coefficient
with the actual residue pullback of $M$.
Every element outside $P$ maps to a nonzero element of $A/P$ and, by
injectivity, of $F$. Both scalar maps therefore extend to $A_P$.
The actual fibre comparison with free localized homology transfers
vanishing from $A/P$ to $F$. The whole canonical base-change isomorphism
identifies the resulting pullback of $M$ with the stated canonical twist
on the same actual $F$-fibre.
\end{proof}

\begin{lemma}[The original canonical coefficient on a prescribed field square]
\label{lem:relative-original-canonical-square-cohomology}
\uses{thm:kn-arithmetic-canonical-prescribed-field-square,lem:kn-scheme-iso-positive-cohomology}
\lean{EnvelopeProperty.Vanishing.RelativeCanonicalSpecialization.squareCanonicalTensor_subsingleton_H_iff}
\leanok
Let $Y/A$ be smooth of relative dimension $n$, and let $L$ be invertible.
Consider a Cartesian field square with original maps $a:X\to Y$ and
$f:X\to\operatorname{Spec}k$. Assume that $X$ has a projective map to
some affine scheme $\operatorname{Spec}R$. Positive cohomology of
$\omega_{X/k}\otimes a^*L$ vanishes exactly when positive cohomology of
the pullback of $\omega_{Y/A}\otimes L$ to the native $k$-fibre vanishes.
There is no properness assumption on $Y/A$ or $X/k$.
\end{lemma}
\begin{proof}
The Cartesian square supplies the scheme isomorphism from the original
$X$ to the native fibre. Its first projection is $a$ and its field
structure is $f$. The whole canonical base-change isomorphism therefore
identifies the pulled coefficient with the stated original canonical
twist. Composing the inverse scheme isomorphism with the given projective
map $X\to\operatorname{Spec}R$ gives a projective map from the native
fibre to that same affine scheme, since an isomorphism is a closed
immersion. The native scheme-isomorphism cohomology comparison applies
using this projective map. The canonical module isomorphism finishes the
comparison with the original twist.
\end{proof}
